\input{../tex-inputs/latex-leseansicht-vorspann}

\section[Gustav Schwarzkopf an Arthur Schnitzler, 27. 2. 1904]{L04314 Gustav Schwarzkopf an Arthur Schnitzler, 27. 2. 1904}
\nopagebreak\mylabel{L04314v}
\rehead{ }\normalsize\beginnumbering\briefempfaengerindex{Schnitzler, Arthur@\textsc{Schnitzler, Arthur}!zzzSchwarzkopf, Gustav@\emph{von Gustav Schwarzkopf}!1904-02-272@{27. 2. 1904}|(be}
\toendnotes[C]{\smallbreak\pagebreak[2]}
\correspDesc{Versand  durch Gustav Schwarzkopf am 27. 2. 1904 in Wien
\newline{}Übermittlung  am 27. 2. 1904 in Wien
\newline{}Zustellung  am 27. 2. 1904 in Wien
\newline{}Erhalt  durch Arthur Schnitzler am 27. 2. 1904 in Wien}\toendnotes[C]{\smallbreak}
\Standort{CUL, Schnitzler, B 96a.}
\physDesc{Postkarte, 414 Zeichen
\newline{}Handschrift: schwarze Tinte, deutsche Kurrent
\newline{}Versand: 1) Rohrpost  2) Stempel: »\nobreak{}\oindex{Wien@\textbf{Wien}, \emph{Verwaltungsgebiet}|pwk}Wien, 27 II 04, 5\textsuperscript{10}N\nobreak{}«.  3) Stempel: »\nobreak{}\oindex{XVIII., Währing@\textbf{XVIII., Währing}, \emph{Verwaltungsgebiet}|pwk}Wien 18, 2\textcolor{gray}{7} II 04, 5\textsuperscript{50}\textcolor{gray}{N}\nobreak{}«. }\toendnotes[C]{\smallbreak}\pstart{}{\pb}Herrn \textsc{D\textsuperscript{r} Arthur Schnitzler}\pend{}\pstart{}XVIII Spöttelgaſſe 7 \textsc{Wien}\oindex{Wien@\textbf{Wien}!XVIII., Währing@\textbf{XVIII., Währing}!Edmund-Weiß-Gasse 7@\textbf{Edmund-Weiß-Gasse 7}, \emph{Wohngebäude}|pw}\pend{}{\bigskip}\vspace{1em}
\pstart
           \raggedleft{}{\pb}27. II. 04\pend
           \vspace{0.5em}
\pstart
           Lieber Arthur! Beſten Dank, aber es geht leider nicht. Max\pwindex{Schwarzkopf, Max 12.\,6.\,1857 Wien – 14.\,4.\,1928 ebd.@\textsc{Schwarzkopf, Max} (12.\,6.\,1857 Wien – 14.\,4.\,1928 ebd.), \emph{Rechtsanwalt}|pw} hat{ }ſchon früher eine Verabredung mit einem Collegen
               getroffen und ich muß ein Böſendorferſaal\oindex{Wien@\textbf{Wien}!I., Innere Stadt@\textbf{I., Innere Stadt}!Bösendorfer-Saal@\textbf{Bösendorfer-Saal}, \emph{Veranstaltungsgebäude}|pw} einen Liederabend\eventindex{Bösendorfer-Saal@\textbf{Bösendorfer-Saal}!Liederabend Karoline Pollatschek, 28.2.1904@Liederabend Karoline Pollatschek, 28.2.1904|pwv} mitmachen, zu dem man mir – \label{K_L04314-1v}\edtext{ich weiß nicht
               warum}{\lemma{\textnormal{\emph{ich weiß nicht
               warum}}}\Cendnote{\textnormal{Als Ansatzpunkt für eine Erklärung könnte man
                  auf die Beziehung zwischen der Sängerin Karoline Pollatschek\pwindex{Pollatschek, Karoline *~4.\,4.\,1867 Wien@\textsc{Pollatschek, Karoline} (*~4.\,4.\,1867 Wien), \emph{Sängerin, Sprachlehrerin}|pwk}
                  und Leo Ebermann\pwindex{Ebermann, Leo 16.\,7.\,1863 Draganovka – 9.\,10.\,1914 Wien@\textsc{Ebermann, Leo} (16.\,7.\,1863 Draganovka – 9.\,10.\,1914 Wien), \emph{Schriftsteller, Journalist, Rechtswissenschaftler}|pwk} verweisen (vgl. A. S.: \emph{Tagebuch}, 7. 12. 1896), mit
                  dem Schwarzkopf\pwindex{Schwarzkopf, Gustav 7.\,11.\,1853 Wien – 13.\,11.\,1939 ebd.@\textsc{Schwarzkopf, Gustav} (7.\,11.\,1853 Wien – 13.\,11.\,1939 ebd.), \emph{Schriftsteller}|pwk} befreundet war.}}}\label{K_L04314-1} – eine Karte in’s Haus geſchickt hat. Da ich nun einmal dankend zugeſagt habe,
               muß ich wol gehen. – Ihre Frau\pwindex{Schnitzler, Olga 17.\,1.\,1882 Wien – 13.\,1.\,1970 Lugano@\textsc{Schnitzler, Olga} (17.\,1.\,1882 Wien – 13.\,1.\,1970 Lugano), \emph{Schauspielerin, Sängerin}|pwv} und Sie herzlichſt grüßend\pend
           
\pstart
           Ihr{\\[\baselineskip]}\spacefill\mbox{Gustav.}\pend
           \leftskip=0em{}\selectlanguage{ngerman}\endnumbering\briefempfaengerindex{Schnitzler, Arthur@\textsc{Schnitzler, Arthur}!zzzSchwarzkopf, Gustav@\emph{von Gustav Schwarzkopf}!1904-02-272@{27. 2. 1904}|)be}\mylabel{L04314h}
\begin{anhang}
\end{anhang}\newcommand{\dateiname}{L04314}\newcommand{\titel}{Gustav Schwarzkopf an Arthur Schnitzler, 27. 2. 1904}\newcommand{\editorInnen}{Herausgegeben von Jahnke, SelmaMüller, Martin Anton}\input{../tex-inputs/latex-leseansicht-abspann}
